% 前言
本书讨论平面上的圆锥曲线问题。与通常的解析几何处理方式不同，本书不把平面上的点写成坐标 $(x,y)$
以后再作代数运算，而是把点与直线都表示为三维实矢量空间中的元素，并直接用矢量与张量的运算
（缩并、点乘、叉乘、对称积、二次叉乘等）表达几何关系。这样做的目的是使几何对象保持为一个
整体，从而使命题的几何内容在代数写法中直接可见。

书中假定读者了解线性代数的基本内容（矢量空间、线性映射、矩阵、二次型、特征值与对角化），
并了解点乘、叉乘与混合积的定义与初等性质。除此之外的内容都在书中给出。

\section*{约定与记号}
\addcontentsline{toc}{section}{约定与记号}

\begin{enumerate}[label=\textbf{(\arabic*)},leftmargin=2.2em]

\item \textbf{指标与求和。} 上指标记矢量，下指标记对偶矢量。一个指标在同一项中上下各出现一次
时表示对该指标从 $1$ 到 $n$ 求和（爱因斯坦求和约定），称为哑指标；其余指标称为自由指标。本书
讨论的是三维矢量空间及其对偶空间，除特别声明外 $i,j,k,l,m,n$ 等指标取 $1,2,3$，二维子空间上
的指标 $\alpha,\beta,\gamma,\delta$ 及其带点形式 $\dot\alpha,\dot\beta$ 等取 $1,2$。$\delta^i_j$ 是恒等张量的分量，
$\delta_{ij}$ 只在标准正交坐标系中有 $1,1,1$ 形式的对角分量，两者不同。

\item \textbf{点矢与线矢。} 小写字母 $a,b,c,p,q,s$ 记矢量（代表平面上的点，称点矢），
$k,l,m,n,u,v$ 记对偶矢量（代表平面上的直线，称线矢），$f,g,h$ 记二阶张量，$A,B,C$ 记对应的
几何对象。矢量与对偶矢量的自然缩并写为 $a\cdot l=a^il_i$。

\item \textbf{度规与升降指标。} 一个对称可逆的 $(0,2)$ 型张量 $g_{ij}$ 称为（矢量空间的）度规，
其逆记为 $(g^{-1})^{ij}$；当讨论中只有一条圆锥曲线并把它取作度规时，简写为 $g^{ij}$，并用它
升降指标：$a_i=g_{ij}a^j$，$l^i=g^{ij}l_j$。

\item \textbf{点乘、叉乘与混合积。} 有了度规之后定义
\[
a\cdot b=g_{ij}a^ib^j,\qquad
(A\times B)_i=\epsilon_{ijk}A^jB^k,\qquad
[A\,B\,C]=\epsilon_{ijk}A^iB^jC^k .
\]
$\epsilon_{ijk}$ 是取定 $\epsilon_{123}=1$ 的三阶反对称张量，$\epsilon^{ijk}$ 由 $\epsilon^{ijk}\epsilon_{ijk}=6$ 确定。

\item \textbf{不写指标的写法。} 若干个张量相邻并列表示指标直接并列而不缩并；两元素之间加点表示
缩并相邻的一个指标；两元素之间加冒号表示缩并相邻的两个指标。例如 $ga$ 的分量为 $g_{ij}a^k$，
$g\cdot a$ 的分量为 $g_{ij}a^j$，而 $f:g=f^{ij}g_{ij}$ 是实数。特别地，两个点（或两条线）的
\textbf{对称积}写为
\[
(a b)^{ij}=a^ib^j+a^jb^i,\qquad
(m n)_{ij}=m_in_j+m_jn_i .
\]

\item \textbf{规范与几何意义。} 三维矢量到平面点的对应是多对一的，故用 $a$ 与 $\lambda a$
（$\lambda\neq0$）表示同一个点。若一个等式在任一元素作这样的替换后仍成立，称它是规范不变的；
规范不变的等式与命题才描述平面射影几何的对象与关系。

\item \textbf{几何层次。} 只用到点矢、线矢及其与坐标系无关的运算的几何称为射影几何；在射影几何上
附加无穷远结构 $\linf$ 后称为仿射几何；再附加正交结构 $\theta$ 后称为正交几何（欧氏几何）；最后
固定坐标系即得解析几何。层次越高可用的运算越多，但表达式越复杂；解题时应在能说清问题的
最低层次中运算。

\item \textbf{二次曲线。} 对称的 $(2,0)$ 型张量 $f^{ij}$ 称为二次点，对称的 $(0,2)$ 型张量 $h_{ij}$
称为二次线，二者统称二次曲线。它们之间的二次点乘与二次叉乘分别为
\[
f:h=f^{ij}h_{ij},\qquad
(f\times h)_{kl}=\epsilon_{imk}\epsilon_{jnl}f^{ij}h^{mn},\qquad
(f\times h)^{kl}=\epsilon^{imk}\epsilon^{jnl}f_{ij}h_{mn},
\]
其中前者的结果是对称 $(0,2)$ 型张量，后者是对称 $(2,0)$ 型张量。不可逆的二次曲线称为退化的。

\item \textbf{编号。} 定义、定理、命题、例、注、习题统一按章连续编号；公式按章编号。

\end{enumerate}

\section*{关于本书的取材}
\addcontentsline{toc}{section}{关于本书的取材}

本书是在一组以三维矢量方法处理圆锥曲线的讲稿基础上重写而成的。重写时删去了原稿中关于
写作缘由与个人感受的叙述，删去了全部插图（图形内容改由文字或由坐标计算给出），并把原稿中
若干处含混或错误的陈述作了更正。原稿中标注为“写得不满意、暂不发表”的第十七节未收入；
原稿第八部分（例题）标注为“已过时”，亦未收入；原稿中关于立体几何与二次曲面的两节
（第二十一、二十二节）属于另一题材，本书不再涉及。
